\section{Related Work}

Our work bridges three research areas: parameter-efficient fine-tuning, sub-quadratic architectures, and transfer learning. We position our contribution at their intersection, revealing architectural constraints that prior work in each area has not systematically explored.

\subsection{Parameter-Efficient Fine-Tuning}

Low-rank adaptation (LoRA) \cite{hu2021lora} introduced a principled approach to fine-tuning large models by injecting low-rank weight updates into linear layers, reducing trainable parameters by orders of magnitude while matching full fine-tuning performance. Subsequent work explored alternative PEFT strategies: prefix tuning \cite{li2021prefix} prepends learnable tokens to input sequences, prompt tuning \cite{lester2021power} optimizes continuous prompts, and adapters \cite{houlsby2019parameter} insert lightweight bottleneck layers. These methods share a common assumption: the base model architecture already possesses the capabilities needed for the target task.

Critically, all established PEFT research focuses exclusively on transformer architectures---BERT \cite{devlin2019bert}, GPT-2 \cite{radford2019language}, T5 \cite{raffel2020exploring}, and their variants. Transformers' bidirectional self-attention enables flexible information flow suitable for diverse tasks. This architectural versatility has allowed PEFT literature to treat the base model as a black box, focusing on \textit{which parameters to adapt} rather than \textit{whether adaptation is possible}. Our work challenges this assumption by testing PEFT transfer to architecturally distinct model families.

\subsection{State-Space Models and Sub-Quadratic Architectures}

State-space models emerged as efficient alternatives to transformers, replacing quadratic attention with linear-time recurrent dynamics. Structured state-space models (S4) \cite{gu2022efficiently} introduced parameterization techniques enabling stable training of deep SSM layers, achieving competitive performance on long-range benchmarks. Mamba \cite{gu2023mamba} extended S4 with selective state-space mechanisms and hardware-aware implementations, demonstrating strong results on language modeling and generation tasks.

However, SSM research has primarily evaluated generative capabilities---language modeling perplexity, long-context understanding, and sequence generation. Classification tasks, especially those requiring bidirectional reasoning, remain underexplored. Mamba's design inherits causal constraints from language modeling: the model processes sequences left-to-right, with each token's representation depending only on preceding context. This architectural choice optimizes generation but may impose fundamental limitations on tasks requiring symmetric comparison.

Other sub-quadratic architectures---linear attention \cite{katharopoulos2020transformers}, RWKV \cite{peng2023rwkv}, and Hyena \cite{poli2023hyena}---share similar motivations but differ in implementation. Our focus on Mamba as a representative causal SSM provides a testbed for understanding how architectural constraints affect PEFT applicability.

\subsection{Transfer Learning and Architectural Compatibility}

Classical transfer learning research \cite{pan2010survey,yosinski2014transferable} established that pretrained representations accelerate downstream task learning. Recent work has explored domain shift \cite{ganin2016domain}, task similarity \cite{zamir2018taskonomy}, and few-shot adaptation \cite{brown2020language}. Yet these studies assume architectural compatibility: the source and target tasks must be solvable by the same model architecture.

Zero-shot evaluation has been used primarily to assess pretrained model quality \cite{radford2019language} rather than as an architectural compatibility test. Our work reframes zero-shot performance as a critical gating mechanism: below-random accuracy signals architectural mismatch that PEFT cannot overcome. This perspective shifts zero-shot evaluation from a benchmark metric to a prerequisite validation step.

\subsection{Positioning Our Contribution}

Prior PEFT research demonstrates \textit{how} to efficiently adapt transformers but does not address \textit{when} adaptation is possible across architectures. SSM research validates efficiency but has not systematically tested architectural limitations on classification tasks. Transfer learning assumes compatibility but lacks principled methods to validate it a priori.

We contribute the first systematic evaluation of PEFT transfer across architecture families, revealing that causal SSMs fail bidirectional tasks at zero-shot evaluation---a failure pattern that LoRA fine-tuning cannot fix. Our task-architecture compatibility framework establishes architectural validation as a prerequisite for PEFT, preventing wasted compute on structurally impossible experiments. As the field explores diverse efficient architectures, understanding these constraints becomes essential for practitioners deploying PEFT methods beyond transformers.
